\chapter{Spherical admissibility and endpoint winding}
\label{chap:sphere_winding}

% Split from the original Gluck_Sphere.tex (stage sections preserved verbatim;
% all labels and \lean pins unchanged).

\paragraph{Formalization map.}
This chapter corresponds primarily to \texttt{Gluck/Sphere/Admissible.lean}, \texttt{Gluck/Sphere/Margins.lean}, \texttt{Gluck/Sphere/FirstVariation.lean}, \texttt{Gluck/Sphere/ConjWinding.lean} and \texttt{Gluck/Sphere/EndpointWinding.lean}.

\section{Admissibility and truncation removal (S2-C)}
\label{sec:sphere_admissible}

\begin{lemma}\leanok
[curvature sensitivity of the truncated speed]
  \label{lem:truncated_speed_sub_le}
  \lean{Gluck.truncatedSpeed_sub_le}
  \uses{def:truncated_speed}
  Let $0\le R$, $0<\delta$ and set $M=\tfrac{1+R^2}{2\delta^2}$. Then for all
  $\kappa,\kappa^\ast:\mathbb R\to\mathbb R$, $\theta$ and $z$,
  \[
    \bigl|\hat q_{\kappa,R,\delta}(\theta,z)-\hat q_{\kappa^\ast,R,\delta}(\theta,z)\bigr|
    \le M\,\bigl|\kappa(\theta)-\kappa^\ast(\theta)\bigr| .
  \]
\end{lemma}

\begin{proof}

\leanok
  The two speeds share the numerator $n=1+(\min(\|z\|,R))^2\in[1,1+R^2]$ and have
  denominators $D=2\max(\kappa(\theta)-c,\delta)$, $D^\ast=2\max(\kappa^\ast(\theta)-c,
  \delta)$ with $c=\langle z,ie^{i\theta}\rangle$; both are $\ge2\delta$, and since
  $x\mapsto\max(x,\delta)$ is $1$-Lipschitz, $|D-D^\ast|\le2|\kappa(\theta)-
  \kappa^\ast(\theta)|$. Hence
  \[
    \Bigl|\frac nD-\frac n{D^\ast}\Bigr|
    =\frac{n\,|D-D^\ast|}{D\,D^\ast}
    \le\frac{(1+R^2)\cdot2|\kappa(\theta)-\kappa^\ast(\theta)|}{4\delta^2}
    =M\,|\kappa(\theta)-\kappa^\ast(\theta)| .
  \]
  (The quotient-difference bound already extracted for
  \cref{lem:truncated_speed_lipschitz} applies verbatim with the numerator difference
  set to zero.)
\end{proof}

\begin{lemma}\leanok
[combined field sensitivity]
  \label{lem:truncated_field_sub_le}
  \lean{Gluck.truncatedField_sub_le}
  \uses{def:truncated_field, lem:truncated_field_lipschitz, lem:truncated_speed_sub_le}
  Let $0\le R$, $0<\delta$, $M=\tfrac{1+R^2}{2\delta^2}$, and let $L$ be \emph{any}
  Lipschitz constant of the truncated field in $z$ (a witness of
  \cref{lem:truncated_field_lipschitz}, consumed as a hypothesis so that one fixed $L$
  can be threaded through the Gr\"onwall argument). Then for all $\kappa,\kappa^\ast$,
  $\theta$, $z$, $\tilde z$,
  \[
    \|F_{\kappa,R,\delta}(\theta,z)-F_{\kappa^\ast,R,\delta}(\theta,\tilde z)\|
    \le L\,\|z-\tilde z\|+M\,|\kappa(\theta)-\kappa^\ast(\theta)| .
  \]
\end{lemma}

\begin{proof}

\leanok
  Insert the mixed term $F_{\kappa,R,\delta}(\theta,\tilde z)$ and use the triangle
  inequality: the first difference is controlled by
  \cref{lem:truncated_field_lipschitz}; the second is
  $(\hat q_{\kappa}(\theta,\tilde z)-\hat q_{\kappa^\ast}(\theta,\tilde z))\,e^{i\theta}$,
  of norm $\le M|\kappa(\theta)-\kappa^\ast(\theta)|$ by
  \cref{lem:truncated_speed_sub_le}.
\end{proof}

\begin{lemma}\leanok
[Gr\"onwall with $L^1$ drive]
  \label{lem:gronwall_L1_drive}
  \lean{Gluck.gronwall_L1_drive}
  Let $T\ge0$, $L\ge0$, $d_0\ge0$, and let $d,g:[0,T]\to\mathbb R$ be continuous with
  $d\ge0$ and $g\ge0$ pointwise. If
  \[
    d(t)\;\le\;d_0+\int_0^t\bigl(L\,d(s)+g(s)\bigr)\,ds
    \qquad\text{for all }t\in[0,T],
  \]
  then, pointwise,
  \[
    d(t)\;\le\;e^{LT}\Bigl(d_0+\int_0^T g\Bigr)\qquad\text{for all }t\in[0,T] .
  \]
  (Project-side infrastructure: Mathlib's Gr\"onwall lemmas bound the derivative by
  $K\|f\|+\varepsilon$ with a \emph{constant} $\varepsilon$; here the drive $g$ is only
  small in $L^1$, which is exactly the regime of the Dahlberg-style reparametrization.
  The Lean statement quantifies the conclusion pointwise over $t\in[0,T]$ rather than
  through a $\sup$; the hypothesis $d\ge0$ is kept for blueprint fidelity though the
  proof does not use it.)
\end{lemma}

\begin{proof}

\leanok
  Set $u(t)=d_0+\int_0^t(L\,d+g)$. The integrand is continuous, so by the fundamental
  theorem of calculus $u$ is differentiable with $u'(t)=L\,d(t)+g(t)$, and $d\le u$ by
  hypothesis, so $u'(t)\le L\,u(t)+g(t)$ ($L\ge0$, and $u\ge d\ge0$). Consider
  $v(t)=e^{-Lt}u(t)-\int_0^t g$: it is differentiable with
  \[
    v'(t)=e^{-Lt}\bigl(u'(t)-L\,u(t)\bigr)-g(t)\le e^{-Lt}g(t)-g(t)\le0,
  \]
  using $0<e^{-Lt}\le1$ and $g\ge0$. A differentiable function with nonpositive
  derivative on an interval is antitone, so $v(t)\le v(0)=d_0$, i.e.
  $e^{-Lt}u(t)\le d_0+\int_0^t g\le d_0+\int_0^T g$ ($g\ge0$ again). Hence
  $d(t)\le u(t)\le e^{Lt}\bigl(d_0+\int_0^T g\bigr)\le e^{LT}\bigl(d_0+\int_0^T g\bigr)$.
\end{proof}

\begin{lemma}

\leanok
[invariant admissible domain — perturbative margin transport]
  \label{lem:invariant_admissible_domain}
  \lean{Gluck.invariant_admissible_domain}
  \uses{def:truncated_speed, def:spherical_flow, lem:spherical_flow_spec,
        lem:truncated_field_sub_le, lem:gronwall_L1_drive,
        lem:sphere_curvature_lower_bound}
  Let $\kappa,\kappa^\ast$ be continuous with $\kappa\ge\kappa_0>0$ pointwise, let
  $0\le R$, $0<\delta$, and let $z^\ast:[0,2\pi]\to\mathbb C$ be a trajectory of the
  $\kappa^\ast$-truncated flow that is \emph{admissible with margin} $\mu>0$:
  \[
    \sup_\theta\|z^\ast(\theta)\|\le R-\mu,
    \qquad
    \sup_\theta\,\langle z^\ast(\theta),\,i e^{i\theta}\rangle\le\kappa_0-\delta-\mu .
  \]
  Let $L$ be any Lipschitz constant of the truncated field in $z$ (a witness of
  \cref{lem:truncated_field_lipschitz}; in Lean $L:\mathbb R_{\ge0}$, coerced to
  $\mathbb R$ inside the exponential) and
  $M=\tfrac{1+R^2}{2\delta^2}$. Then every trajectory $z$ of the $\kappa$-truncated flow
  with
  \[
    e^{2\pi L}\Bigl(\|z(0)-z^\ast(0)\|+M\!\int_0^{2\pi}\!|\kappa-\kappa^\ast|\Bigr)\le\mu
  \]
  is admissible: $\|z(\theta)\|\le R$ and
  $\kappa(\theta)-\langle z(\theta),i e^{i\theta}\rangle\ge\delta$ for all
  $\theta\in[0,2\pi]$.
\end{lemma}

\begin{proof}

\leanok
  \emph{Field sensitivity.} By \cref{lem:truncated_field_sub_le},
  \[
    \|F_{\kappa}(\theta,z)-F_{\kappa^\ast}(\theta,\tilde z)\|
    \le L\|z-\tilde z\|+M\,|\kappa(\theta)-\kappa^\ast(\theta)| .
  \]
  \emph{Integral inequality.} Both trajectories are $C^1$ on $[0,2\pi]$ with derivatives
  given by their fields (\cref{lem:spherical_flow_spec}), so by the fundamental theorem
  of calculus $z(\theta)-z^\ast(\theta)=z(0)-z^\ast(0)+\int_0^\theta\bigl(F_\kappa(s,
  z(s))-F_{\kappa^\ast}(s,z^\ast(s))\bigr)ds$; taking norms and using the field
  sensitivity, $d(\theta)=\|z(\theta)-z^\ast(\theta)\|$ satisfies
  $d(\theta)\le d(0)+\int_0^\theta\bigl(L\,d+M|\kappa-\kappa^\ast|\bigr)$, with $d$
  continuous and the drive $g=M|\kappa-\kappa^\ast|$ continuous and nonnegative.
  \emph{Gr\"onwall with $L^1$ drive.} \cref{lem:gronwall_L1_drive} gives
  $\sup_{[0,2\pi]}d\le e^{2\pi L}\bigl(d(0)+M\|\kappa-\kappa^\ast\|_{L^1[0,2\pi]}\bigr)
   \le\mu$.
  \emph{Margin propagation.} Then
  $\|z(\theta)\|\le\|z^\ast(\theta)\|+\mu\le R$, and since
  $|\langle z-z^\ast,\,i e^{i\theta}\rangle|\le d(\theta)\le\mu$,
  \[
    \kappa(\theta)-\langle z,i e^{i\theta}\rangle
    \ge\kappa_0-\langle z^\ast,i e^{i\theta}\rangle-\mu\ge\delta .
  \]
  \emph{Why perturbative is the right shape.} No direct confinement estimate uniform in
  $\kappa_0$ exists: the radial velocity $d\|z\|^2/d\theta=2q\langle z,e^{i\theta}\rangle$
  changes sign, no disk is flow-invariant (at a radial critical point with
  $c=\langle z,ie^{i\theta}\rangle=-\|z\|$ one computes that a local maximum of $\|z\|$ is
  allowed exactly when $\kappa(\theta)\ge(1-\|z\|^2)/(2\|z\|)$), and for
  $\kappa_0<1/\sqrt3$ the disk of the model radius already contains denominator-vanishing
  states — admissibility is irreducibly a condition on the pair
  $(\|z\|,\langle z,ie^{i\theta}\rangle)$. The margins are instead supplied by an
  \emph{explicit} model (the symmetric step bicircle of \cref{sec:sphere_winding}, whose
  arcs give $\langle z^\ast,ie^{i\theta}\rangle=\langle w_j,ie^{i\theta}\rangle-r_j
  \le\|w_j\|-r_j$ on the $j$-th arc), and the $L^1$ distance is made arbitrarily small by
  the Dahlberg-style reparametrization of \cref{chap:dahlberg_step1}.
\end{proof}

\begin{definition}\leanok
[periodic extension of a fundamental-window function]
  \label{def:periodic_extension}
  \lean{Gluck.periodicExtension}
  For $z:\mathbb R\to\mathbb C$, the \emph{$2\pi$-periodic extension} is
  $\tilde z(t)=z\bigl(t-2\pi\lfloor t/(2\pi)\rfloor\bigr)$: every $t$ is reduced by an
  integer multiple of $2\pi$ into the fundamental window $[0,2\pi)$ and $z$ is evaluated
  there. This is the device promised in \cref{lem:reconstruction_ode} for promoting a
  trajectory defined on $[0,2\pi]$ to a globally defined closed curve.
\end{definition}

\begin{lemma}\leanok
[the periodic extension is periodic]
  \label{lem:periodic_extension_periodic}
  \lean{Gluck.periodicExtension_periodic}
  \uses{def:periodic_extension}
  For every $z$, the extension $\tilde z$ of \cref{def:periodic_extension} is
  $2\pi$-periodic: $\tilde z(t+2\pi)=\tilde z(t)$ for all $t$.
\end{lemma}

\begin{proof}

\leanok
  $\lfloor(t+2\pi)/(2\pi)\rfloor=\lfloor t/(2\pi)\rfloor+1$, so the reduced arguments
  agree: $t+2\pi-2\pi(\lfloor t/(2\pi)\rfloor+1)=t-2\pi\lfloor t/(2\pi)\rfloor$.
\end{proof}

\begin{lemma}

\leanok
[truncation removal: admissible closed trajectories realize $\kappa$]
  \label{lem:reconstruction_ode}
  \lean{Gluck.reconstruction_ode}
  \uses{lem:truncated_speed_eq, lem:spherical_flow_spec, lem:truncated_speed_pos,
        lem:spherical_speed_solve, def:realizes_spherical_curvature,
        lem:spherical_speed_periodic, def:periodic_extension,
        lem:periodic_extension_periodic}
  Let $\kappa$ be continuous and $2\pi$-periodic (positivity NOT assumed — see the
  re-sign note at the end of this statement), $0<\delta$, $R<1$, and let
  $z:[0,2\pi]\to\mathbb C$ be a trajectory of the truncated flow that is admissible
  ($\|z(\theta)\|\le R$ and $\kappa(\theta)-\langle z(\theta),ie^{i\theta}\rangle\ge\delta$
  for all $\theta$) and closed ($z(2\pi)=z(0)$). Then the $2\pi$-periodic extension
  $\tilde z$ of $z$ (\cref{def:periodic_extension}) agrees with $z$ on $[0,2\pi]$, is a
  closed curve, and realizes $\kappa$ (\cref{def:realizes_spherical_curvature}) in the
  gauge $\varphi(\theta)=\theta$ — in particular it is $C^1$, confined to the open unit
  disk, and solves the \emph{true} reconstruction ODE $z'=q_\kappa(\theta,z)e^{i\theta}$.
  (The Lean statement packages the conclusion as the conjunction
  $\mathrm{EqOn}\wedge\mathrm{IsClosedCurve}\wedge\mathrm{RealizesSphericalCurvature}$
  for $\tilde z$, \emph{plus a fourth conjunct exposing the true-ODE derivative
  globally}: $\forall t,\ \tilde z$ has at $t$ the derivative
  $q_\kappa(t,\tilde z(t))\,e^{it}$. The proof already establishes this fact on the way
  to the realization predicate; exposing it is what lets the simplicity chain
  (\cref{lem:spherical_trajectory_eq_reconstruct}) identify $\tilde z$ with a translated
  Euclidean reconstruction curve without re-deriving the seam-matching argument. The
  hypothesis $0\le R$ turned out to be unnecessary and is dropped. Additionally, the
  re-sign RELAXES the hypothesis ``$\kappa$ is a curvature function'' to ``$\kappa$ is
  continuous and $2\pi$-periodic'' — positivity of $\kappa$ enters the proof only
  through the admissibility margin $\kappa-\langle z,ie^{i\theta}\rangle\ge\delta$,
  never through $\kappa>0$ itself, and the relaxed form is exactly what stage 2
  (\cref{chap:sphere_mixed}) needs, so one declaration serves both stages. Fallback if
  the relaxation surfaces a genuine positivity use: keep the positive hypothesis and
  record the obstacle in the task result; the derivative conjunct alone still unblocks
  the S2-E chain.)
\end{lemma}

\begin{proof}

\leanok
  On the admissible set $\hat q=q_\kappa$ (\cref{lem:truncated_speed_eq}), so the flow
  equation of \cref{lem:spherical_flow_spec} becomes the true ODE on $[0,2\pi]$; the
  right-hand side is continuous, so $z$ is $C^1$ there. Closedness plus
  $2\pi$-periodicity of the field ($\kappa$ periodic,
  \cref{lem:spherical_speed_periodic}) make the periodic extension a global $C^1$
  solution (one-sided derivatives match at the seam because the field values match).
  The speed $q_\kappa>0$ along the trajectory (\cref{lem:spherical_speed_solve} shape;
  positivity from admissibility), so $z'\ne0$, the tangent angle is $\varphi(\theta)=
  \theta$ itself, $\|z'\|=q_\kappa$, and the speed relation of
  \cref{def:realizes_spherical_curvature} is exactly the algebraic solve of
  \cref{def:spherical_speed}. Confinement $\|z\|\le R<1$ gives the open-disk condition.
\end{proof}

\section{Endpoint winding over the symmetric step model (S2-D)}
\label{sec:sphere_winding}

\begin{lemma}\leanok
[arc bracket identity]
  \label{lem:constant_arc_inner}
  \lean{Gluck.constant_arc_inner}
  For all $w\in\mathbb C$, $r,\theta\in\mathbb R$:
  $\bigl\langle w-i\,r\,e^{i\theta},\,i e^{i\theta}\bigr\rangle
   =\langle w,ie^{i\theta}\rangle-r$.
\end{lemma}

\begin{proof}

\leanok
  Bilinearity of the real inner product and
  $\langle ie^{i\theta},ie^{i\theta}\rangle=\|ie^{i\theta}\|^2=1$
  (\texttt{real\_inner\_smul\_left}, \texttt{real\_inner\_self\_eq\_norm\_sq}).
\end{proof}

\begin{lemma}\leanok
[arc norm expansion]
  \label{lem:constant_arc_norm_sq}
  \lean{Gluck.constant_arc_norm_sq}
  For all $w\in\mathbb C$, $r,\theta\in\mathbb R$:
  $\|w-i\,r\,e^{i\theta}\|^2=\|w\|^2-2r\langle w,ie^{i\theta}\rangle+r^2$.
\end{lemma}

\begin{proof}

\leanok
  The real polarization expansion $\|u-v\|^2=\|u\|^2-2\langle u,v\rangle+\|v\|^2$
  (\texttt{norm\_sub\_sq\_real}) with $\|i\,r\,e^{i\theta}\|^2=r^2$ and
  \cref{lem:constant_arc_inner}-style bilinearity for the cross term.
\end{proof}

\begin{lemma}\leanok
[consistency identity of the arc data]
  \label{lem:constant_arc_consistency}
  \lean{Gluck.constant_arc_consistency}
  \uses{def:spherical_speed, lem:constant_arc_norm_sq, lem:constant_arc_inner}
  Fix $K$ and start data $(\theta_0,z_0)$ with positive bracket
  $K-\langle z_0,ie^{i\theta_0}\rangle>0$, and set $r=q_K(\theta_0,z_0)$,
  $w=z_0+i\,r\,e^{i\theta_0}$. Then $1+\|w\|^2=2rK+r^2$.
\end{lemma}

\begin{proof}

\leanok
  Expand $\|w\|^2=\|z_0\|^2+2r\langle z_0,ie^{i\theta_0}\rangle+r^2$
  (\texttt{norm\_add\_sq\_real}), then substitute the defining relation
  $2r\bigl(K-\langle z_0,ie^{i\theta_0}\rangle\bigr)=1+\|z_0\|^2$ of
  $r=q_K(\theta_0,z_0)$ (\cref{def:spherical_speed}).
\end{proof}

\begin{lemma}

\leanok
[constant-curvature arcs are explicit circular arcs]
  \label{lem:constant_curvature_arc}
  \lean{Gluck.constant_curvature_arc}
  \uses{def:spherical_speed, lem:spherical_flow_unique, lem:constant_arc_inner,
        lem:constant_arc_norm_sq, lem:constant_arc_consistency}
  Fix $K>0$, a start $(\theta_0,z_0)$ with $K-\langle z_0,ie^{i\theta_0}\rangle>0$, and set
  \[
    r=\frac{1+\|z_0\|^2}{2\bigl(K-\langle z_0,ie^{i\theta_0}\rangle\bigr)}=q_K(\theta_0,z_0),
    \qquad w=z_0+i\,r\,e^{i\theta_0}.
  \]
  Then $z(\theta)=w-i\,r\,e^{i\theta}$ solves the true reconstruction ODE
  $z'=q_K(\theta,z)e^{i\theta}$ with constant curvature $K$, pointwise at every $\theta$
  where the bracket stays positive ($K-\langle z(\theta),ie^{i\theta}\rangle>0$): at each
  such $\theta$ one has $q_K(\theta,z(\theta))=r$ and hence
  $z'(\theta)=q_K(\theta,z(\theta))\,e^{i\theta}$. Along the arc
  $\langle z(\theta),ie^{i\theta}\rangle=\langle w,ie^{i\theta}\rangle-r$. Moreover the
  Euclidean data satisfy the \emph{consistency identity} $K=(1+\|w\|^2-r^2)/(2r)$
  (generalizing the centered case $w=0$ of \cref{lem:spherical_speed_circle}).
\end{lemma}

\begin{proof}

\leanok
  Differentiate: $z'(\theta)=r e^{i\theta}$, so the claim is $q_K(\theta,z(\theta))=r$.
  Compute $\langle z,ie^{i\theta}\rangle=\langle w,ie^{i\theta}\rangle-r$ and
  $\|z\|^2=\|w\|^2-2r\langle w,ie^{i\theta}\rangle+r^2$; then
  $1+\|z\|^2=2r(K-\langle z,ie^{i\theta}\rangle)$ reduces to
  $1+\|w\|^2=2rK+r^2$, i.e.\ the consistency identity, which holds at $\theta_0$ by the
  definitions of $r$ and $w$ (expand $\|w\|^2=\|z_0\|^2+2r\langle z_0,ie^{i\theta_0}
  \rangle+r^2$). Spherically: stereographic images of spherical circles are Euclidean
  circles, and a constant-$K$ arc through $(\theta_0,z_0)$ is the circle of Euclidean
  radius $r=q$ tangent to $e^{i\theta_0}$ at $z_0$. Identification with the flow (where
  admissible, with clamps inactive) is \cref{lem:spherical_flow_unique}.
\end{proof}

\begin{lemma}\leanok
[half-turn invariance of the truncated speed]
  \label{lem:truncated_speed_half_turn}
  \lean{Gluck.truncatedSpeed_half_turn}
  \uses{def:truncated_speed}
  If $\kappa(\theta+\pi)=\kappa(\theta)$, then
  $\hat q_{\kappa,R,\delta}(\theta+\pi,-z)=\hat q_{\kappa,R,\delta}(\theta,z)$ for all
  $\theta,z$.
\end{lemma}

\begin{proof}

\leanok
  $\|{-z}\|=\|z\|$ and $ie^{i(\theta+\pi)}=-ie^{i\theta}$, so
  $\langle -z,ie^{i(\theta+\pi)}\rangle=\langle z,ie^{i\theta}\rangle$; with
  $\kappa(\theta+\pi)=\kappa(\theta)$ every ingredient of the clamped quotient is
  unchanged.
\end{proof}

\begin{lemma}\leanok
[half-turn anti-equivariance of the truncated field]
  \label{lem:truncated_field_half_turn}
  \lean{Gluck.truncatedField_half_turn}
  \uses{def:truncated_field, lem:truncated_speed_half_turn}
  If $\kappa(\theta+\pi)=\kappa(\theta)$, then
  $F_{\kappa,R,\delta}(\theta+\pi,-z)=-F_{\kappa,R,\delta}(\theta,z)$ for all
  $\theta,z$.
\end{lemma}

\begin{proof}

\leanok
  $F(\theta+\pi,-z)=\hat q(\theta+\pi,-z)\,e^{i(\theta+\pi)}
   =\hat q(\theta,z)\cdot(-e^{i\theta})=-F(\theta,z)$ by
  \cref{lem:truncated_speed_half_turn} and $e^{i\pi}=-1$.
\end{proof}

\begin{lemma}

\leanok
[half-turn equivariance]
  \label{lem:flow_half_turn_equivariance}
  \lean{Gluck.flow_half_turn_equivariance}
  \uses{def:truncated_speed, lem:truncated_field_solution_unique, lem:spherical_flow_spec,
        lem:truncated_speed_half_turn, lem:truncated_field_half_turn}
  If $\kappa(\theta+\pi)=\kappa(\theta)$ for all $\theta$, then the truncated field is
  invariant under the half turn: $\hat q(\theta+\pi,-z)=\hat q(\theta,z)$, hence if
  $z(\cdot)$ solves the truncated ODE on $[0,2\pi]$ then $y(\theta)=-z(\theta+\pi)$
  solves it on $[0,\pi]$. Consequently, writing $H$ for the half-period map
  $z_0\mapsto\Phi(z_0,\pi)$ of a $\pi$-periodic curvature: if $H(z_0)=-z_0$ then the
  trajectory through $z_0$ satisfies the central symmetry $z(\theta+\pi)=-z(\theta)$ for
  $\theta\in[0,\pi]$; in particular it closes at $2\pi$: $z(2\pi)=-z(\pi)=z_0$.
\end{lemma}

\begin{proof}

\leanok
  $\|{-z}\|=\|z\|$ and $\langle -z,\,ie^{i(\theta+\pi)}\rangle=\langle -z,\,-ie^{i\theta}
  \rangle=\langle z,ie^{i\theta}\rangle$, and $\kappa(\theta+\pi)=\kappa(\theta)$; so
  every ingredient of $\hat q$ is unchanged. For $\theta\in[0,\pi]$ one has
  $\theta+\pi\in[\pi,2\pi]\subseteq[0,2\pi]$, so
  $y'(\theta)=-z'(\theta+\pi)=-\hat q\,e^{i(\theta+\pi)}=\hat q(\theta,y(\theta))e^{i\theta}$
  holds on $[0,\pi]$ — $y$ solves the truncated ODE there (only there: this is why the
  comparison below happens on the half interval). If $H(z_0)=-z_0$, then
  $y(0)=-z(\pi)=z_0=z(0)$, and $z$ restricted to $[0,\pi]$ solves the same ODE
  (\cref{lem:spherical_flow_spec}), so $y=z$ on $[0,\pi]$ by the two-solution
  uniqueness on $[0,\pi]$ (\cref{lem:truncated_field_solution_unique} with $T=\pi$),
  i.e.\ $z(\theta+\pi)=-z(\theta)$ for $\theta\in[0,\pi]$; at $\theta=\pi$ this gives
  $z(2\pi)=-z(\pi)=z_0$.
\end{proof}

\subsection*{Design overview (S2-D, resolved iter-048)}

The two formerly-open sub-questions of this phase are now resolved at design level, and
the whole phase is elementary algebra plus the in-tree winding toolkit. The shape:

\textbf{DO-NOT (degeneracy of the constant model).} For \emph{constant} $\kappa=K$ every
admissible initial point lies on a spherical circle of curvature $K$
(\cref{lem:constant_curvature_arc}), so every trajectory closes and the endpoint error
map vanishes identically: the constant family carries no degree and must NOT be the
homotopy target. The fix is the \emph{symmetric step bicircle} with \emph{small
contrast}: levels $a=c-h/2$, $b=c+h/2$ around the midpoint $c$ of the four-vertex
overlap window, on the in-tree canonical equal-quarter breakpoints $0,\pi/2,\pi,3\pi/2$
(\cref{def:four_arc_step_function} — note $\kappa^\ast=$
$\texttt{stepCurvature } b\,a\,0\,(\pi/2)\,\pi\,(3\pi/2)$ takes the value $a$ on
$[0,\pi/2)\cup[\pi,3\pi/2)$ and $b$ on the complement, and is $\pi$-periodic).

\textbf{Everything is explicit arc algebra.} Because the step is piecewise constant, its
model trajectories are concatenations of \emph{explicit circular arcs}
(\cref{lem:constant_curvature_arc}); no flow, no variational equation, and no
$\mathrm{PicardLindel\ddot of}$ package is needed on the model side. The model endpoint
error map $E^\ast_{a,b}$ (\cref{def:step_error_map}) is a closed-form composition of
four arc maps. Its degeneracy at $h=0$ becomes a feature: $E^\ast_{c,c}\equiv0$
identically, so $E^\ast_{a,b}=h\,W+O(h^2)$ for an explicit \emph{first-variation field}
$W$, and the planner's symbolic computation (numerically validated: the arc composite
matches an RK4 integration of the step ODE to $10^{-7}$, $E^\ast(z_0^\ast)=O(h^2)$, and
$E^\ast/h$ matches the predicted linearization with ratio $\to1$) gives
\[
  W(z_0^\ast)=0,\qquad
  DW(z_0^\ast)=-\tfrac{2r^\ast}{1+c^2}\cdot\overline{(\,\cdot\,)}
  \qquad(z_0^\ast=-i\,r^\ast,\ r^\ast=\sqrt{1+c^2}-c),
\]
a \emph{conjugation-type} (orientation-reversing) nondegenerate linear map: the model
error loop on a small circle $\partial B(z_0^\ast,\rho)$ has winding $-1\ne0$. The exact
mechanism behind the cancellations is the \emph{quadratic identity}
\cref{lem:speed_quadratic_identity}: $q_c(\theta,z)-r^\ast$ vanishes to second order in
the distance to the centered circle, with an exact rational formula (no Taylor
remainders anywhere: all expansions in this section are exact algebraic identities plus
norm estimates of explicitly-controlled factors).

\textbf{Transport.} The comparison of the actual (reparametrized) $\kappa$-flow with the
step model cannot use \cref{lem:invariant_admissible_domain} directly ($\kappa^\ast$ is
discontinuous, and the concatenated model has corners), and Mathlib-style continuity of
the drive fails at the breakpoints. The resolution is \emph{arcwise chaining}: on each
quarter arc the model level is a \emph{constant} $K\in\{a,b\}$, the drive
$M|\kappa-K|$ is continuous there, and the single-arc transport
(\cref{lem:invariant_admissible_arc}) is the verbatim margin-transport proof on a
shifted interval. Four chained applications give the full-period transport
(\cref{lem:step_model_transport}) with constant $e^{2\pi L}$, both for admissibility and
for the endpoint estimate $\|E_\kappa(z_0)-E^\ast(z_0)\|\le e^{2\pi L}M\!\int|\kappa-
\kappa^\ast|$.

\textbf{Quantifier order.} $c$ and $\kappa_0=\min\kappa$ are fixed by $\kappa$; margins
$(R,\delta,\mu,\rho_0,h_0)$ by \cref{lem:step_model_margins}; the disk radius $\rho$ and
contrast $h$ by \cref{lem:step_error_expansion} (boundary margin
$\sim\tfrac{r^\ast}{1+c^2}h\rho$); the reparametrization quality $\varepsilon\ll h\rho$
last (it is free, \cref{lem:step_L1_reparam}). The Euclidean
\texttt{errorMap\_winding\_eq\_one} \emph{value} is not reusable; only
\cref{thm:existence_of_zero} and the perturbation/congruence lemmas of
\cref{chap:winding} are.

\begin{definition}\leanok
[spherical arc map]
  \label{def:spherical_arc_map}
  \lean{Gluck.sphericalArcMap}
  \uses{def:spherical_speed}
  For $K,\theta_0,\Delta\in\mathbb R$ and $z\in\mathbb C$, the \emph{arc map} is
  \[
    A_{K,\theta_0,\Delta}(z)\;=\;z+i\,q_K(\theta_0,z)\,e^{i\theta_0}\bigl(1-e^{i\Delta}\bigr),
  \]
  where $q_K$ abbreviates the gauge speed $q_{(\theta\mapsto K)}$ of
  \cref{def:spherical_speed} for the constant curvature $K$. By
  \cref{lem:constant_curvature_arc}, where the bracket stays positive this is the
  time-$\Delta$ endpoint of the constant-$K$ trajectory started at $(\theta_0,z)$: the
  arc trajectory is $\theta\mapsto w-i\,r\,e^{i\theta}$ with $r=q_K(\theta_0,z)$,
  $w=z+i\,r\,e^{i\theta_0}$, and its value at $\theta_0+\Delta$ is exactly
  $A_{K,\theta_0,\Delta}(z)$. (A total function of the junk-value kind, like $q$ itself:
  the geometric meaning requires bracket positivity, the definition does not.)
\end{definition}

\begin{lemma}\leanok
[quadratic identity: exact second-order vanishing of $q$ at the centered circle]
  \label{lem:speed_quadratic_identity}
  \lean{Gluck.sphericalSpeed_sub_radius}
  \uses{def:spherical_speed, lem:spherical_speed_circle}
  Fix $c>0$ and set $r^\ast=\sqrt{1+c^2}-c$, so that $r^{\ast2}+2cr^\ast-1=0$. For all
  $\theta$ and all $z$ with $D:=c-\langle z,ie^{i\theta}\rangle\ne0$,
  \[
    q_c(\theta,z)-r^\ast\;=\;\frac{\bigl\|z+r^\ast i e^{i\theta}\bigr\|^2}{2D}
    \;=\;\frac{\|z-z_c(\theta)\|^2}{2D},
    \qquad z_c(\theta):=-r^\ast i e^{i\theta}
  \]
  ($z_c$ is the centered circle of \cref{lem:spherical_speed_circle}). In particular
  $q_c-r^\ast$ vanishes \emph{exactly to second order} in the distance to the centered
  circle: the $z$-gradient of $q_c(\theta,\cdot)$ vanishes at $z_c(\theta)$, and on
  $D>0$ one has $q_c\ge r^\ast$.
\end{lemma}

\begin{proof}
  \uses{lem:constant_arc_norm_sq, lem:constant_arc_inner}
  \leanok
  $q_c-r^\ast=\dfrac{(1+\|z\|^2)-2r^\ast D}{2D}$, and the numerator is
  \[
    1+\|z\|^2-2r^\ast c+2r^\ast\langle z,ie^{i\theta}\rangle
    =\|z\|^2+r^{\ast2}+2\langle z,\,r^\ast ie^{i\theta}\rangle
    =\|z+r^\ast ie^{i\theta}\|^2,
  \]
  using $1-2r^\ast c=r^{\ast2}$ (the defining identity of $r^\ast$) and the polarization
  expansion (\cref{lem:constant_arc_norm_sq} with $-r$ replaced by $r^\ast$).
\end{proof}

\begin{lemma}\leanok
[speed dominates the centered radius]
  \label{lem:speed_radius_le}
  \lean{Gluck.sphericalSpeed_radius_le}
  \uses{def:spherical_speed, lem:speed_quadratic_identity}
  Fix $c>0$ and $r^\ast=\sqrt{1+c^2}-c$. Wherever
  $D=c-\langle z,ie^{i\theta}\rangle>0$ one has $r^\ast\le q_c(\theta,z)$ — the
  inequality half of \cref{lem:speed_quadratic_identity}, used to keep model arcs
  outside the centered circle.
\end{lemma}

\begin{proof}
  \uses{lem:speed_quadratic_identity}
  \leanok
  Immediate from \cref{lem:speed_quadratic_identity}: $q_c-r^\ast$ is a square divided
  by $2D>0$, hence nonnegative.
\end{proof}

\begin{lemma}\leanok
[exact level sensitivity of the gauge speed]
  \label{lem:speed_level_sensitivity}
  \lean{Gluck.sphericalSpeed_sub_level}
  \uses{def:spherical_speed}
  For constant levels $K,K'$ and any $(\theta,z)$ with both brackets
  $D_K=K-\langle z,ie^{i\theta}\rangle$ and $D_{K'}=K'-\langle z,ie^{i\theta}\rangle$
  nonzero,
  \[
    q_K(\theta,z)-q_{K'}(\theta,z)=\frac{(1+\|z\|^2)\,(K'-K)}{2\,D_K\,D_{K'}} .
  \]
  This is mechanism (i) of \cref{lem:step_error_expansion}: the level sensitivity of
  the gauge speed is an exact one-line quotient identity — no Taylor remainder.
\end{lemma}

\begin{proof}

\leanok
  Both speeds are the quotients $(1+\|z\|^2)/(2D_K)$ and $(1+\|z\|^2)/(2D_{K'})$
  (\cref{def:spherical_speed} at constant curvature); subtract and clear denominators.
\end{proof}

\begin{lemma}\leanok
[speed half-turn (true speed)]
  \label{lem:spherical_speed_half_turn}
  \lean{Gluck.sphericalSpeed_half_turn}
  \uses{def:spherical_speed}
  For every constant $K$: $q_K(\theta+\pi,-z)=q_K(\theta,z)$ for all $\theta,z$. (Same
  computation as \cref{lem:truncated_speed_half_turn}, without clamps: $\|-z\|=\|z\|$ and
  $\langle -z,ie^{i(\theta+\pi)}\rangle=\langle z,ie^{i\theta}\rangle$.)
\end{lemma}

\begin{proof}
  \uses{lem:truncated_speed_half_turn}
  \leanok
  As stated; $ie^{i(\theta+\pi)}=-ie^{i\theta}$ and bilinearity of the inner product.
\end{proof}

\begin{lemma}\leanok
[arc-map half-turn anti-equivariance]
  \label{lem:arc_map_half_turn}
  \lean{Gluck.sphericalArcMap_half_turn}
  \uses{def:spherical_arc_map, lem:spherical_speed_half_turn}
  $A_{K,\theta_0+\pi,\Delta}(-z)=-A_{K,\theta_0,\Delta}(z)$ for all
  $K,\theta_0,\Delta,z$.
\end{lemma}

\begin{proof}

\leanok
  $q_K(\theta_0+\pi,-z)=q_K(\theta_0,z)$ (\cref{lem:spherical_speed_half_turn}) and
  $e^{i(\theta_0+\pi)}=-e^{i\theta_0}$, so every term of
  \cref{def:spherical_arc_map} flips sign.
\end{proof}

\begin{lemma}\leanok
[arc concatenation]
  \label{lem:arc_map_concat}
  \lean{Gluck.sphericalArcMap_concat}
  \uses{def:spherical_arc_map, lem:constant_curvature_arc}
  Fix $K>0$, start data $(\theta_0,z)$ with positive bracket, and $\Delta_1,\Delta_2\ge0$.
  If the bracket $K-\langle z(\theta),ie^{i\theta}\rangle$ stays positive along the arc
  trajectory $z(\theta)=w-ire^{i\theta}$ for $\theta\in[\theta_0,\theta_0+\Delta_1+
  \Delta_2]$, then
  \[
    A_{K,\theta_0+\Delta_1,\Delta_2}\bigl(A_{K,\theta_0,\Delta_1}(z)\bigr)
    =A_{K,\theta_0,\Delta_1+\Delta_2}(z) .
  \]
\end{lemma}

\begin{proof}

\leanok
  By \cref{lem:constant_curvature_arc}(iii), $q_K$ is constant $=r$ along the arc where
  the bracket is positive; in particular
  $q_K(\theta_0+\Delta_1,\,z(\theta_0+\Delta_1))=r=q_K(\theta_0,z)$, so the second arc
  continues the same circle: both sides equal $w-ire^{i(\theta_0+\Delta_1+\Delta_2)}$
  after expanding \cref{def:spherical_arc_map} and collecting the exponentials.
  In particular the full-turn map is the identity where admissible:
  $A_{K,\theta_0,2\pi}(z)=z$ (since $e^{2\pi i}=1$), which is the exact form of the
  constant-model degeneracy.
\end{proof}

\begin{definition}\leanok
[step half-period map]
  \label{def:step_half_map}
  \lean{Gluck.stepHalfMap}
  \uses{def:spherical_arc_map}
  For levels $a,b$, the \emph{half-period map} of the symmetric equal-quarter step is
  the composition of the two quarter-arc maps
  \[
    H_{a,b}(z)\;=\;A_{b,\pi/2,\pi/2}\bigl(A_{a,0,\pi/2}(z)\bigr) .
  \]
  (First quarter has level $a$, second $b$, matching
  $\kappa^\ast=\texttt{stepCurvature } b\,a\,0\,(\pi/2)\,\pi\,(3\pi/2)$ of
  \cref{def:four_arc_step_function}.)
\end{definition}

\begin{definition}\leanok
[step model error map]
  \label{def:step_error_map}
  \lean{Gluck.stepErrorMap}
  \uses{def:step_half_map}
  $E^\ast_{a,b}(z)\;=\;-H_{a,b}\bigl(-H_{a,b}(z)\bigr)-z$. By
  \cref{lem:arc_map_half_turn}, the arcs of the second half period (levels $a,b$ at
  $\theta_0=\pi,3\pi/2$, by $\pi$-periodicity of the step) are obtained from the first
  half by the half turn, so $z+E^\ast_{a,b}(z)$ is exactly the four-arc composite
  \[
    A_{b,3\pi/2,\pi/2}\circ A_{a,\pi,\pi/2}\circ A_{b,\pi/2,\pi/2}\circ A_{a,0,\pi/2}
  \]
  — the endpoint at $2\pi$ of the concatenated step-model trajectory. At $a=b=c$,
  $E^\ast_{c,c}(z)=0$ wherever the circle through $z$ stays bracket-positive
  (\cref{lem:arc_map_concat}).
\end{definition}

\begin{lemma}\leanok
[four-arc composite form of the step error map]
  \label{lem:step_error_map_four_arc}
  \lean{Gluck.stepErrorMap_four_arc}
  \uses{def:step_error_map, def:step_half_map, lem:arc_map_half_turn}
  For all levels $a,b$ and all $z$,
  \[
    z+E^\ast_{a,b}(z)
    =A_{b,3\pi/2,\pi/2}\Bigl(A_{a,\pi,\pi/2}\bigl(A_{b,\pi/2,\pi/2}(A_{a,0,\pi/2}(z))
    \bigr)\Bigr) :
  \]
  the reflected double-half-period form of \cref{def:step_error_map} coincides with
  the four-arc composite — the endpoint at $2\pi$ of the concatenated step-model
  trajectory.
\end{lemma}

\begin{proof}
  \uses{lem:arc_map_half_turn}
  \leanok
  \cref{lem:arc_map_half_turn} at $\theta_0=0$ and $\theta_0=\pi/2$ rewrites the two
  arcs of the second half period: $A_{a,\pi,\pi/2}(y)=-A_{a,0,\pi/2}(-y)$ and
  $A_{b,3\pi/2,\pi/2}(y)=-A_{b,\pi/2,\pi/2}(-y)$. Unfolding
  $E^\ast_{a,b}(z)=-H_{a,b}(-H_{a,b}(z))-z$ and cancelling the double negations gives
  the composite.
\end{proof}

\begin{lemma}\leanok
[explicit arcs solve the truncated ODE]
  \label{lem:constant_arc_solves_truncated}
  \lean{Gluck.constant_arc_solves_truncated}
  \uses{lem:constant_curvature_arc, lem:truncated_speed_eq, def:truncated_field}
  Fix a constant $K$, $0<\delta$, $R$, start data $(\theta_0,z_0)$, and let
  $z(\theta)=w-ire^{i\theta}$ be the arc trajectory of \cref{lem:constant_curvature_arc}.
  If along $\theta\in[t_1,t_2]$ the arc is \emph{admissible with clamps inactive} —
  $\|z(\theta)\|\le R$ and $K-\langle z(\theta),ie^{i\theta}\rangle\ge\delta$ — then $z$
  solves the \emph{truncated} ODE for the constant curvature $K$ on $[t_1,t_2]$ in the
  $\mathrm{HasDerivWithinAt}$ sense:
  $z'(\theta)=\hat q_{K,R,\delta}(\theta,z(\theta))\,e^{i\theta}$ on $[t_1,t_2]$.
\end{lemma}

\begin{proof}

\leanok
  On the admissible set the clamps are inactive, $\hat q=q_K$
  (\cref{lem:truncated_speed_eq}), and $z'(\theta)=re^{i\theta}
  =q_K(\theta,z(\theta))e^{i\theta}$ by \cref{lem:constant_curvature_arc}.
\end{proof}

\begin{lemma}\leanok
[single-arc margin transport (shifted interval, constant model level)]
  \label{lem:invariant_admissible_arc}
  \lean{Gluck.invariant_admissible_arc}
  \uses{lem:truncated_field_sub_le, lem:gronwall_L1_drive,
        lem:invariant_admissible_domain}
  Let $\kappa$ be continuous with $\kappa\ge\kappa_0$ pointwise, $0<\delta$, let $L$ be a
  Lipschitz witness of \cref{lem:truncated_field_lipschitz} and
  $M=\tfrac{1+R^2}{2\delta^2}$. Let $[t_1,t_2]\subseteq\mathbb R$, $K\in\mathbb R$
  (\emph{constant} model level), and let $z,z^\ast:[t_1,t_2]\to\mathbb C$ be continuous
  with, in the $\mathrm{HasDerivWithinAt}$ sense on $[t_1,t_2]$,
  \[
    z'=\hat q_{\kappa,R,\delta}(\theta,z)\,e^{i\theta},\qquad
    z^{\ast\prime}=\hat q_{K,R,\delta}(\theta,z^\ast)\,e^{i\theta},
  \]
  and suppose the model has margins $\mu>0$ on the arc:
  $\|z^\ast(\theta)\|\le R-\mu$ and
  $\langle z^\ast(\theta),ie^{i\theta}\rangle\le\kappa_0-\delta-\mu$ for
  $\theta\in[t_1,t_2]$. If
  \[
    e^{L(t_2-t_1)}\Bigl(\|z(t_1)-z^\ast(t_1)\|
      +M\!\int_{t_1}^{t_2}\!|\kappa(s)-K|\,ds\Bigr)\;\le\;\mu,
  \]
  then for all $\theta\in[t_1,t_2]$:
  \[
    \|z(\theta)-z^\ast(\theta)\|\le
    e^{L(t_2-t_1)}\Bigl(\|z(t_1)-z^\ast(t_1)\|+M\!\int_{t_1}^{t_2}\!|\kappa-K|\Bigr),
    \quad
    \|z(\theta)\|\le R,\quad
    \kappa(\theta)-\langle z(\theta),ie^{i\theta}\rangle\ge\delta .
  \]
\end{lemma}

\begin{proof}

\leanok
  Verbatim the proof of \cref{lem:invariant_admissible_domain} after the time shift
  $s\mapsto t_1+s$ (which turns $[t_1,t_2]$ into $[0,T]$, $T=t_2-t_1$, and shifts the
  fields accordingly): FTC difference identity, the field sensitivity
  \cref{lem:truncated_field_sub_le} with $\kappa^\ast\equiv K$ — the drive
  $g=M|\kappa-K|$ is \emph{continuous} because the model level is constant, which is the
  entire point of the arcwise formulation — then \cref{lem:gronwall_L1_drive} and margin
  propagation via $|\langle z-z^\ast,ie^{i\theta}\rangle|\le\|z-z^\ast\|$.
\end{proof}

\begin{definition}\leanok
[quarter-arc margin package]
  \label{def:arc_margins}
  \lean{Gluck.arcMargins}
  \uses{def:spherical_speed, def:spherical_arc_map}
  For data $(\kappa_0,R,\delta,\mu,K,t_1,t_2)$ and a start point $p$, write
  $\mathrm{arcMargins}(\kappa_0,R,\delta,\mu,K,t_1,t_2,p)$ for the statement that
  along the constant-level-$K$ arc trajectory
  $z^\ast(\theta)=w-i\,r\,e^{i\theta}$ through $(t_1,p)$ (with $r=q_K(t_1,p)$,
  $w=p+i\,r\,e^{it_1}$), for every $\theta\in[t_1,t_2]$:
  \[
    \|z^\ast(\theta)\|\le R-\mu,\qquad
    \langle z^\ast(\theta),ie^{i\theta}\rangle\le\kappa_0-\delta-\mu,\qquad
    K-\langle z^\ast(\theta),ie^{i\theta}\rangle\ge\delta .
  \]
  This packages, for one quarter arc, exactly the model-margin hypotheses of
  \cref{lem:invariant_admissible_arc} together with the clamp-inactivity hypothesis
  of \cref{lem:constant_arc_solves_truncated}.
\end{definition}

\begin{lemma}\leanok
[one quarter-arc of the step transport]
  \label{lem:quarter_step_transport}
  \lean{Gluck.quarter_step_transport}
  \uses{def:arc_margins, def:spherical_arc_map, def:truncated_field,
        lem:truncated_field_lipschitz}
  Let $\kappa$ be continuous with $\kappa\ge\kappa_0$, $0\le R$, $0<\delta$,
  $t_1\le t_2$, $L$ a Lipschitz witness for the truncated field, and
  $M=\tfrac{1+R^2}{2\delta^2}$. Let $z$ solve the $\kappa$-truncated ODE on
  $[t_1,t_2]$ (in the $\mathrm{HasDerivWithinAt}$ sense), let $p\in\mathbb C$ carry
  $\mathrm{arcMargins}(\kappa_0,R,\delta,\mu,K,t_1,t_2,p)$, and suppose
  \[
    e^{L(t_2-t_1)}\Bigl(\|z(t_1)-p\|+M\!\int_{t_1}^{t_2}\!|\kappa-K|\Bigr)\le\mu .
  \]
  Then $z$ is admissible on $[t_1,t_2]$ ($\|z\|\le R$, bracket $\ge\delta$) and
  \[
    \bigl\|z(t_2)-A_{K,t_1,t_2-t_1}(p)\bigr\|
    \le e^{L(t_2-t_1)}\Bigl(\|z(t_1)-p\|+M\!\int_{t_1}^{t_2}\!|\kappa-K|\Bigr)
  \]
  — the single step of the \cref{lem:step_model_transport} chain.
\end{lemma}

\begin{proof}
  \uses{lem:invariant_admissible_arc, lem:constant_arc_solves_truncated,
        lem:constant_arc_consistency}
        \leanok
  The clamp-inactivity component of \cref{def:arc_margins} at $\theta=t_1$ gives
  bracket positivity at the start, hence the consistency identity
  (\cref{lem:constant_arc_consistency}) for the arc data $(w,r)$; by
  \cref{lem:constant_arc_solves_truncated} the model arc solves the truncated ODE on
  the quarter (its norm is $\le R-\mu\le R$ and its level-$K$ bracket $\ge\delta$ by
  the package). \cref{lem:invariant_admissible_arc} then transports the distance
  bound and admissibility; the model arc endpoint at $t_2$ is exactly
  $A_{K,t_1,t_2-t_1}(p)$ by \cref{def:spherical_arc_map}.
\end{proof}

\begin{lemma}\leanok
[four-arc chained transport against the step model]
  \label{lem:step_model_transport}
  \lean{Gluck.stepModel_transport}
  \uses{def:arc_margins, def:spherical_arc_map, lem:invariant_admissible_arc,
        lem:constant_arc_solves_truncated,
        def:step_half_map, def:step_error_map, lem:arc_map_half_turn,
        def:four_arc_step_function, def:spherical_flow, lem:spherical_flow_spec}
  Let $\kappa$ be continuous with $\kappa\ge\kappa_0$ pointwise, $0<\delta$, $L$, $M$ as
  above, levels $a,b$, and $\kappa^\ast=\texttt{stepCurvature } b\,a\,0\,(\pi/2)\,\pi\,
  (3\pi/2)$. Let $z$ be a trajectory of the $\kappa$-truncated flow on $[0,2\pi]$ with
  $z(0)=z_0$, and let the \emph{concatenated step model} from $z_0$ be the four
  quarter-arc trajectories with levels $a,b,a,b$ at $\theta_0=0,\pi/2,\pi,3\pi/2$, each
  arc started at the endpoint of the previous one. Assume each of the four arcs
  carries the margin package
  $\mathrm{arcMargins}(\kappa_0,R,\delta,\mu,K_j,\theta_j,\theta_{j+1},p_j)$ of
  \cref{def:arc_margins}, where $\theta_j=j\pi/2$, $K_j$ alternates $a,b,a,b$, and
  $p_j$ is the $j$-fold nested arc-map start point ($p_0=z_0$,
  $p_{j+1}=A_{K_j,\theta_j,\pi/2}(p_j)$), and
  \[
    e^{2\pi L}\,M\!\int_0^{2\pi}\!|\kappa-\kappa^\ast|\;\le\;\mu .
  \]
  Then $z$ is admissible on $[0,2\pi]$ ($\|z\|\le R$, bracket $\ge\delta$), and the
  endpoint satisfies
  \[
    \bigl\|\,\bigl(z(2\pi)-z_0\bigr)-E^\ast_{a,b}(z_0)\,\bigr\|
    \;\le\;e^{2\pi L}\,M\!\int_0^{2\pi}\!|\kappa-\kappa^\ast| .
  \]
\end{lemma}

\begin{proof}
  \uses{lem:quarter_step_transport, lem:step_error_map_four_arc}
  \leanok
  Apply \cref{lem:quarter_step_transport} on each quarter $[\theta_j,\theta_{j+1}]$,
  $\theta_j=j\pi/2$, with the constant level $K_j\in\{a,b\}$ of that quarter; on each
  quarter $\int_{\theta_j}^{\theta_{j+1}}|\kappa-K_j|=\int_{\theta_j}^{\theta_{j+1}}
  |\kappa-\kappa^\ast|$ ($\kappa^\ast=K_j$ off the single breakpoint). Chain the distance
  bounds: with $d_j=\|z(\theta_j)-p_j\|$ and $I_j=M\int_{\theta_j}^{
  \theta_{j+1}}|\kappa-\kappa^\ast|$, each quarter gives $d_{j+1}\le e^{L\pi/2}(d_j+I_j)$
  and $d_0=0$, whence inductively $d(\theta)\le e^{2\pi L}\sum_j I_j$ throughout —
  in particular each quarter's smallness hypothesis holds and admissibility propagates
  across all four quarters. The model endpoint at $2\pi$ is the four-arc composite,
  which equals $z_0+E^\ast_{a,b}(z_0)$ by \cref{lem:step_error_map_four_arc}, giving
  the endpoint estimate.
\end{proof}

\begin{lemma}\leanok
[uniform margins of the step model near the centered circle]
  \label{lem:step_model_margins}
  \lean{Gluck.stepModel_margins}
  \uses{def:arc_margins, def:spherical_arc_map, lem:constant_arc_inner,
        lem:speed_quadratic_identity, lem:speed_radius_le,
        lem:spherical_speed_circle, def:step_half_map}
  Fix $c>0$ and set $r^\ast=\sqrt{1+c^2}-c$, $z_0^\ast=-i\,r^\ast$, and let
  $\kappa_0>-r^\ast$. (Stage-2 re-sign: the lemma landed in stage 1 with the
  hypothesis $0<\kappa_0$; the mixed-sign stage \cref{chap:sphere_mixed} RELAXES it
  to $-r^\ast<\kappa_0$ — same conclusion, same constants ledger. Positivity of
  $\kappa_0$ enters the landed proof \emph{only} through positivity of the slack
  $m=\min(1-r^\ast,\ \kappa_0+r^\ast)$, which holds exactly when
  $\kappa_0>-r^\ast$; the three slack inequalities fed to the per-arc margin
  step use only $m\le\kappa_0+r^\ast$ and $m>0$, never the sign of $\kappa_0$.
  The stage-1 consumer (now the ε-generic \cref{thm:spaceform_endpoint_winding}) supplies
  $-r^\ast<\kappa_0$ trivially from $0<\kappa_0$ and $r^\ast>0$. Fallback if the
  relaxation surfaces a genuine positivity use: keep the positive form and add a
  separate relaxed variant, recording the obstacle in the task result.)
  There exist explicit $0<R<1$ (the positivity witness is needed
  by \cref{lem:step_model_transport}'s hypothesis $0\le R$), $\delta>0$, $\mu>0$,
  $\rho_0>0$, $h_0>0$ (functions of $c,\kappa_0$ only) such that for all levels $a,b$
  with $|a-c|\le h_0$, $|b-c|\le h_0$ and every $z_0$ with $\|z_0-z_0^\ast\|\le\rho_0$,
  the four quarter-arc margin packages of \cref{lem:step_model_transport} hold:
  \[
    \mathrm{arcMargins}(\kappa_0,R,\delta,\mu,K_j,\theta_j,\theta_{j+1},p_j),
    \qquad j=0,1,2,3,
  \]
  with $\theta_j=j\pi/2$, levels $K_j$ alternating $a,b,a,b$, and nested start points
  $p_0=z_0$, $p_{j+1}=A_{K_j,\theta_j,\pi/2}(p_j)$ (\cref{def:arc_margins}; the third
  component keeps the clamps inactive so \cref{lem:constant_arc_solves_truncated}
  applies inside \cref{lem:quarter_step_transport}).
\end{lemma}

\begin{proof}
  \uses{lem:spherical_speed_circle, lem:constant_arc_inner, lem:speed_radius_le}
  \leanok
  Along the centered circle itself ($a=b=c$, $z_0=z_0^\ast$) one has
  $\|z^\ast\|\equiv r^\ast<1$ and $\langle z^\ast(\theta),ie^{i\theta}\rangle\equiv
  -r^\ast<0$ (\cref{lem:spherical_speed_circle}), so all three quantities have strict
  slack: norm slack $1-r^\ast$, bracket slacks $\kappa_0+r^\ast$ and $c+r^\ast$. A
  concrete admissible ledger (the prover may adjust constants, recording its choices
  in the task result): $R=(1+r^\ast)/2$ (norm slack $\tfrac{1-r^\ast}{2}$),
  $\delta=\mu=\tfrac18\min\bigl(1-r^\ast,\ \kappa_0+r^\ast\bigr)$, so that
  $\delta+\mu\le\tfrac{\kappa_0+r^\ast}{4}$ and (since $\kappa_0\le c$) also
  $\le\tfrac{c+r^\ast}{4}$; at the centered circle all three inequalities then hold
  with residual slack $\ge\tfrac12\min\bigl(\tfrac{1-r^\ast}{2},
  \tfrac{\kappa_0+r^\ast}{2}\bigr)=:\sigma$. Stability: each quarter arc is
  $\theta\mapsto w_j-i r_j e^{i\theta}$ with $r_j=q_{K_j}(\theta_j,p_j)$,
  $w_j=p_j+i r_j e^{i\theta_j}$; on the compact neighborhood
  $\{|K-c|\le c/2,\ \|z-z_c(\theta_j)\|\le r^\ast/2\}$ the bracket
  $K-\langle z,ie^{i\theta_j}\rangle$ is bounded below by an explicit positive
  constant (Cauchy–Schwarz plus \cref{lem:constant_arc_inner}; cf.\
  \cref{lem:speed_radius_le} for the lower bound $q\ge r^\ast$ on the
  positive-bracket region), so $q_K(\theta_j,\cdot)$ is a quotient of Lipschitz
  functions with denominator bounded away from $0$ and $(K,p)\mapsto(w,r)$, hence each
  arc map $A_{K,\theta_j,\pi/2}$, is Lipschitz with an explicit constant $C_1(c)$.
  Along each arc,
  $\langle z^\ast(\theta),ie^{i\theta}\rangle=\langle w_j,ie^{i\theta}\rangle-r_j$
  with $(w_j,r_j)$ within $C_1(c)(\rho_0+h_0)$ of $(0,r^\ast)$. Chaining the four
  Lipschitz maps, every point of the concatenated trajectory and every pair
  $(w_j,r_j)$ deviates from the centered-circle data by at most
  $C(c)\,(\rho_0+h_0)$ with $C(c)=C_1(c)^4+C_1(c)^3+C_1(c)^2+C_1(c)+1$ (a crude but
  explicit chaining constant); choose $\rho_0=h_0=\sigma/(4C(c))$. All constants are
  explicit; no compactness argument is needed.
\end{proof}

\begin{lemma}\leanok
[first-variation expansion of the step error map]
  \label{lem:step_error_expansion}
  \lean{Gluck.stepError_expansion}
  \uses{def:step_error_map, def:step_half_map, def:spherical_arc_map,
        lem:speed_quadratic_identity, lem:speed_level_sensitivity,
        lem:step_error_map_four_arc, lem:arc_map_concat, lem:constant_arc_inner}
  Fix $c>0$; let $r^\ast,z_0^\ast$ be as above and set $\eta=\tfrac{2r^\ast}{1+c^2}>0$.
  There exist explicit $\rho_1>0$, $h_1>0$, $C>0$ (functions of $c$ only) such that for
  all $h\in(0,h_1]$, with levels $a=c-h/2$, $b=c+h/2$, and all $z_0$ with
  $\|z_0-z_0^\ast\|\le\rho_1$:
  \[
    \Bigl\|\,E^\ast_{a,b}(z_0)\;+\;\eta\,h\,\overline{(z_0-z_0^\ast)}\,\Bigr\|
    \;\le\;C\,h\,\bigl(\|z_0-z_0^\ast\|^2+h\bigr) .
  \]
  \emph{Provenance.} The linear part $-\eta\,h\,\overline{(\,\cdot\,)}$ is
  planner-derived (iter-048) and numerically validated ($E^\ast/h$ against the predicted
  field at multiple base points and directions, ratio $\to1$ as the offset shrinks);
  the prover should re-derive the coefficient symbolically during formalization and may
  correct its exact value — the \emph{load-bearing} content is only that the linear part
  is a nonzero real multiple of complex conjugation (orientation-reversing,
  nondegenerate).
\end{lemma}

\begin{proof}

\leanok
  All expansions are exact algebra; no abstract Taylor theorem is used. Write
  $\delta z=z_0-z_0^\ast$. The two mechanisms:
  (i) \emph{level sensitivity is exact}: the LANDED lemma
  \cref{lem:speed_level_sensitivity} (\texttt{Gluck.sphericalSpeed\_sub\_level}) gives,
  for any levels $K,K'$ and admissible $z$,
  \[
    q_K(\theta,z)-q_{K'}(\theta,z)
    =\frac{(1+\|z\|^2)\,(K'-K)}{2\,D_K D_{K'}},\qquad
    D_K=K-\langle z,ie^{i\theta}\rangle ;
  \]
  at the circle base point this evaluates to
  $\mp\tfrac{h}{2}\cdot\tfrac{r^\ast}{\sqrt{1+c^2}}+O(h(h+\|\delta z\|))$ for
  $K=c\pm h/2$.
  (ii) \emph{base-point sensitivity is exactly quadratic}:
  \cref{lem:speed_quadratic_identity} gives $q_c(\theta,z)-r^\ast=\|z-z_c(\theta)\|^2/
  (2D)$, so all first-order-in-$\delta z$ contributions at level $c$ vanish and
  re-enter only through the $h$-dependent terms.
  Expand the four-arc composite $z_0+E^\ast_{a,b}(z_0)$
  (\cref{lem:step_error_map_four_arc}) arc by arc using (i), (ii) and
  the gradient formula for $q_c$ implied by (ii) (each arc's output is the input plus
  $i\,q\,e^{i\theta_0}(1-e^{i\pi/2})$ with $1-e^{i\pi/2}=1-i$, and intermediate points
  stay within $O(\|\delta z\|+h)$ of the centered circle by
  \cref{lem:step_model_margins}-style bounds); after cancellation of the $h$-free part
  ($E^\ast_{c,c}\equiv0$, \cref{lem:arc_map_concat}) the coefficient of $h$ collects to
  $-\eta\,\overline{\delta z}+O(\|\delta z\|^2)$ and every remaining term carries
  $h(\|\delta z\|^2+h)$. The bookkeeping is finite (eight $q$-evaluations), and the
  prover should organize it through named intermediate \texttt{have}-identities rather
  than one monolithic \texttt{ring} call.
\end{proof}

\begin{definition}\leanok
[conjugate boundary loop]
  \label{def:conj_loop}
  \lean{Gluck.conjLoop}
  For $w\in\mathbb C$, the \emph{conjugate loop} is the continuous map
  $[0,1]\to\mathbb C$, $t\mapsto w\cdot\overline{e^{2\pi it}}=w\,e^{-2\pi it}$ — the
  once-clockwise circle of radius $|w|$ through $w$. It is the model loop for the
  boundary behaviour of the conjugation-type first variation of
  \cref{lem:step_error_expansion}.
\end{definition}

\begin{lemma}\leanok
[the conjugate loop avoids the origin]
  \label{lem:conj_loop_ne_zero}
  \lean{Gluck.conjLoop_ne_zero}
  \uses{def:conj_loop}
  For $w\ne0$ and every $t\in[0,1]$, $w\cdot\overline{e^{2\pi it}}\ne0$: it is a product
  of two nonzero factors ($|e^{-2\pi it}|=1$).
\end{lemma}

\begin{proof}
  \uses{def:conj_loop}
  \leanok
  Immediate: $|w\,e^{-2\pi it}|=|w|\ne0$.
\end{proof}

\begin{lemma}\leanok
[winding of a conjugate loop]
  \label{lem:winding_conj_loop}
  \lean{Gluck.windingNumberC_conj_loop}
  \uses{def:winding_number_c, def:conj_loop, lem:conj_loop_ne_zero}
  For $w\ne0$, the loop $t\mapsto w\cdot\overline{e^{2\pi it}}$ on $[0,1]$ is nowhere
  zero and has $\mathbb C$-winding number $-1$.
\end{lemma}

\begin{proof}

\leanok
  The normalization of the loop is $t\mapsto\mathrm{circleProj}(w)\cdot e^{-2\pi it}$,
  which lifts through $\mathrm{Circle.exp}$ via $t\mapsto\arg(w)-2\pi t$; by
  the same lift computation as \cref{lem:winding_number_eq_div_of_lift} — which is
  \texttt{private} and must be replicated locally, see the CAUTION below — the winding
  number is
  $\bigl((\arg w-2\pi)-\arg w\bigr)/(2\pi)=-1$. (CAUTION — the \emph{entire} lift layer
  of \texttt{Winding.lean} is \texttt{private} and not importable: not just
  \texttt{windingNumber\_negStandard} but also \texttt{windingNumber\_eq\_div\_of\_lift},
  \texttt{angleLift}, and the raw \texttt{windingNumber}; the public surface is only
  \texttt{windingNumberC}, \texttt{windingNumberC\_congr},
  \texttt{windingNumberC\_eq\_of\_perturb}, \texttt{windingNumberC\_posScalarField},
  \texttt{diskBoundaryLoop}, and \texttt{thm:existence\_of\_zero}. The local proof must
  therefore redo the angle-lift computation self-contained from Mathlib's
  \texttt{Circle.isCoveringMap\_exp} path lifting — the same route as
  \texttt{Winding.lean}'s construction, an estimated 80--120 LOC helper layer, NOT a
  five-line replication.)
\end{proof}

\begin{lemma}\leanok
[four values at freely chosen levels]
  \label{lem:exists_abab_levels}
  \lean{Gluck.exists_abab_levels}
  \uses{def:four_vertex_condition, lem:ivt_hits, lem:abab_values}
  Let $\kappa$ be continuous and $2\pi$-periodic with alternating extrema data
  $p_1<q_1<p_2<q_2<p_1+2\pi$ satisfying the strict value separation
  $\max(\kappa(q_1),\kappa(q_2))<\min(\kappa(p_1),\kappa(p_2))$. Then for \emph{every}
  pair of levels $a,b$ with
  $\max(\kappa(q_1),\kappa(q_2))<a<b<\min(\kappa(p_1),\kappa(p_2))$ there exist
  $\theta_1<\theta_2<\theta_3<\theta_4<\theta_1+2\pi$ with
  $\kappa(\theta_1,\theta_2,\theta_3,\theta_4)=(a,b,a,b)$. (This refines
  \cref{lem:abab_values}, which produces one specific pair of levels — the endpoints of
  the overlap window; here the levels are free, which is what allows the small-contrast
  choice $a=c-h/2$, $b=c+h/2$.)
\end{lemma}

\begin{proof}

\leanok
  Four applications of the packaged IVT (\cref{lem:ivt_hits}): on $[q_1,p_2]$, $\kappa$
  runs from $\le a$ to $\ge b$, giving $\theta_1$ with $\kappa=a$ and then, on
  $[\theta_1,p_2]$, $\theta_2\ge\theta_1$ with $\kappa=b$ (strict inequality since the
  values differ); on $[p_2,q_2]$ ($\ge b\to\le a$) a point $\theta_3$ with $\kappa=a$;
  on $[q_2,p_1+2\pi]$ ($\le a\to\kappa(p_1)\ge b$, using periodicity) a point
  $\theta_4$ with $\kappa=b$. The ordering and window conditions follow since
  $\theta_4\le p_1+2\pi<q_1+2\pi\le\theta_1+2\pi$. (Home: this lemma lives in the S²
  file — the Euclidean files are frozen.)
\end{proof}

\begin{lemma}\leanok
[$L^1$ step reparametrization]
  \label{lem:step_L1_reparam}
  \lean{Gluck.exists_step_L1_reparam}
  \uses{lem:exists_preliminary_reparam, def:four_arc_step_function,
        def:is_curvature_function}
  Let $\kappa$ be a curvature function, $0<a<b$, and $\theta_1<\theta_2<\theta_3<
  \theta_4<\theta_1+2\pi$ with $\kappa(\theta_1,\theta_2,\theta_3,\theta_4)=(a,b,a,b)$.
  For every $\varepsilon>0$ there is an orientation-preserving circle reparametrization
  $h_1$ ($\mathrm{StrictMono}$, $C^1$ with continuous positive derivative,
  $h_1(\theta+2\pi)=h_1(\theta)+2\pi$) such that
  \[
    \int_0^{2\pi}\bigl|\kappa(h_1(\theta))-\kappa^\ast(\theta)\bigr|\,d\theta
    \;<\;\varepsilon,
    \qquad \kappa^\ast=\texttt{stepCurvature } b\,a\,0\,(\pi/2)\,\pi\,(3\pi/2).
  \]
\end{lemma}

\begin{proof}

\leanok
  Apply \cref{lem:exists_preliminary_reparam} with
  $\varepsilon'=\varepsilon/\bigl(B+2\pi+1\bigr)$ where
  $B=\sup_{[0,2\pi]}\kappa+b$ (finite by compactness and periodicity, and a pointwise
  bound for $|\kappa\circ h_1-\kappa^\ast|$ since $0<\kappa^\ast\le b$ and $\kappa\circ
  h_1$ ranges over the values of $\kappa$). Split the integral over the bad set (measure
  $<\varepsilon'$, integrand $\le B$) and its complement in $[0,2\pi)$ (integrand
  $\le\varepsilon'$, measure $\le2\pi$):
  $\int\le\varepsilon' B+2\pi\varepsilon'<\varepsilon$.
\end{proof}


